% \documentclass{report}

% \usepackage{graphicx}
% \usepackage{dcolumn}
% \usepackage{bm}
% \usepackage{amsfonts, amssymb, amsmath}
% \usepackage{braket}


%\begin{document}

\subsection{Localized shallow semiconductor heterostructures in the bulk}
\label{Shallow}
The interruption of the crystalline structure of the semiconductor, as with
impurities or heterostructures, can serve to capture excess electrons within the
semiconductor. We wish to examine the physics of these captured electrons. The
Hamiltonian of such a system may be written as
\begin{equation}
  \left[
    \frac{\hat{\bm{p}}^2}{2 m_e} + \hat{V}^0 + \hat{U} - E
  \right]\ket{\Psi} = 0,
\end{equation}
where $\hat{V}^0$ is the periodic potential of the crystalline lattice. We
introduce $\hat{U}$ as a potential introduced by defects in the lattice such as
that generated by an impurity. $E$ has energy zero at the lowest conduction band
minima.

We understand that, in the absence of defect, the wavefunction resulting from
the periodic crystal potential will be a Bloch function. The acceptable energies
for free electrons occur in discrete bands that may be separated, resulting in
disallowed energies. Each Bloch function is denoted by a wavevector and an index
indicating membership in one of these bands. The space wavefunction of each
Bloch state follows the form
\begin{equation}
\braket{\bm{r} | n\bm{k}}
  =
    \frac{1}{\left(2\pi\right)^{\frac{3}{2}}}
    \mathrm{e}^{\im\bm{k}\cdot\bm{r}}
    u_{n\bm{k}\bm{r}},
\end{equation}
where $u_{n\bm{k}}$ may be a complicated function, but one which always has the
periodicity of the lattice. The Bloch state is the eigenstate of a nearly free
electron traveling through a silicon lattice. As such, the Schr\"{o}dinger
equation of the pure system may be written as
\begin{equation}
\left[\frac{\hat{\bm{p}}^2}{2 m_e} + \hat{V}^0 - E_{n\bm{k}}\right]
\ket{n\bm{k}}
  =
    0
\end{equation}
$\hat{V}^0$ is the periodic potential of silicon and $E_{n\bm{k}}$ is the energy
of a nearly free electron traveling through silicon in the $n$th conduction
band given the wavevector $k$.

We assume that Bloch functions taken from the first Brillouin zone and from all
bands form a complete orthonormal basis to be used to construct the state of an
excess electron bound to an impurity or heterostructure.  As such, we represent
the complete wavefunction as
\begin{equation}
  \ket{\Psi} = \sum_n \int_{BZ} \diff\bm{k} \, A_{n\bm{k}}\ket{n\bm{k}}
\label{definition:bloch-basis:wavefunction}
\end{equation}
with the Bloch states such that
\begin{equation}
  \braket{n\bm{k} | n^{\prime}\bm{k}^{\prime}}
  =
  \delta_{nn^{\prime}}\delta\left(\bm{k}-\bm{k}^{\prime}\right)
\end{equation}
We write the Schr\"{o}dinger equation in the Bloch basis, letting $\hat{U}$ be
the defect potential,
\begin{equation}
\left(E_{n\bm{k}} - E\right)A_{n\bm{k}}
  +
\sum_{n^\prime}\int_{BZ}\diff\bm{k}^\prime\,A_{n^\prime\bm{k}^\prime}
\braket{n\bm{k} | \hat{U} | n^\prime\bm{k}^\prime}
  =
    0
\end{equation}

We can clearly see the role of the defect potential in mixing together different
bands, as well as causing the state to spread through each band. However, this
equation proves difficult to solve directly, since it relies on the interplay
between the Fourier components of the defect potential and the periodic function
$u_{n\bm{k}}$.

To resolve these difficulties we employ a variational method. In the variational
method one calculates $\mathcal{E}\left[\Psi\right]$ from
\begin{equation}
\mathcal{E}\left[\Psi\right]
  = 
    \frac
    { \Bra{\Psi}\frac{\hat{\bm{p}}^2}{2 m_e} + \hat{V}^0 + \hat{U}\Ket{\Psi} }
    { \braket{\Psi | \Psi} },
\label{definition:expectation:energy}
\end{equation}
where the state $\Psi$ is judiciously chosen, with variational
parameters, to solve the problem approximately. Minimization of the
energy with respect to the parameters provides the best possible
approximate solution to the problem using the state
$\Psi$.



\subsection{Shallow one electron states}
At this point, we explicitly restrict our concern to localized ``shallow''
potentials. For our purposes, the ``shallowness'' of a potential depends on its
effect in mixing and occupying the free electron states. Namely, the bound state
of a shallow potential only significantly occupies the lowest conduction band
centered steeply around band minima. As such, the bound state possesses
non-periodic structure that emerges at a scale much longer then the
characteristic wavelength of the Bloch function at the conduction band minima.
The shallowness assumptions leads to a series of conclusions:
\begin{enumerate}
\item The potential does not cause significant band occupation outside the
  lowest conduction band, and so terms $n \ne 0$ may be ignored.
\item The band occupation will mostly occur in ``valleys'' -- the $k$-space
  energy minima of the lowest conduction band, which we will denote as
  $\bm{k}_{q}$. We write the occupation as $A_{0\bm{k}} = \sum_{q}
  A^{(q)}_{\bm{k}}$. Above, the $A^{(q)}_{\bm{k}}$ functions are sharply peaked
  functions centered at the $q$th minimum of the conduction band. We expect that
  the overall form of of these occupation peaks will obey symmetry relations and
  will follow a form corresponding to one of the irreducible representations of
  the symmetry group of the composite Hamiltonian including the $\hat{U}$ term.
  We presume the set containing the $A^{(q)}$ functions is closed under the
  transformations of the relevant symmetry group.
\end{enumerate}
We wish to see the effect of these assumptions upon Eq.
\eqref{definition:expectation:energy}. Since we are considering only the lowest
conduction band, we can drop the explicit band index $n=0$ from our
computations. 

The diagonal part is straightforward, and so we write
\begin{equation}
\mathcal{T}
  \equiv
    \Braket{\Psi | \frac{\hat{\bm{p}}^2}{2 m_e} + \hat{V}^0 | \Psi}
      =
        \int_{BZ}
        \diff\bm{k}\,
        \tilde{A}^{(q)}_{\bm{k}}
        E^{(q^{\prime})}_{\bm{k}}
        A^{(q^{\prime})}_{\bm{k}}.
  \label{element:diagonal}
\end{equation}

Turning our concern to the impurity term in Eq.
\eqref{definition:expectation:energy} yields
\begin{equation}
  \mathcal{U}
  \equiv
  \braket{\Psi | \hat{U} | \Psi}
  =
  \sum_{q,q^\prime}
  \iint_{BZ}\diff\bm{k}\,\diff\bm{k}^\prime\,
  \tilde{A}^{(q)}_{\bm{k}}  A^{(q^\prime)}_{\bm{k}^\prime}
  \braket{\bm{k} | \hat{U} | \bm{k}^\prime},
\end{equation}
where the matrix element is computed as as
\begin{equation}
  \braket{\bm{k} | \hat{U} | \bm{k}^\prime}
  =
  \frac{1}{\left(2\pi\right)^{3}}
  \int \diff\bm{r}\,
  U_{\bm{r}}
  \mathrm{e}^{\im\left(\bm{k}^{\prime} - \bm{k}\right)\cdot\bm{r}}
  \tilde{u}_{\bm{k}\bm{r}}
  u_{\bm{k}^{\prime}\bm{r}}.
\end{equation}
We may express the impurity term as a space integral,
\begin{equation}
  \mathcal{U}
  =
  \int \diff\bm{r}\,
  U_{\bm{r}}
  \tilde{\Psi}^{\left(q^\prime\right)}_{\bm{r}}
  \Psi^{\left(q\right)}_{\bm{r}},
\end{equation}
where $\Psi^{(q)}$ is
\begin{equation}
  \Psi^{(q)}_{\bm{r}}
  =
  \frac{1}{\left(2\pi\right)^{\frac{3}{2}}}
  \int_{BZ}\diff\bm{k}\,
  A^{(q)}_{\bm{k}}
  \mathrm{e}^{\im\bm{k}\cdot\bm{r}}
  u_{\bm{k}\bm{r}}.
\label{definition:bloch-basis:wavefunction:single-valley}
\end{equation}
The significance of the function $\Psi^{(q)}$ as a single-valley term of the
full multi-valley wavefunction is clear when we write the space wavefunction of
the bound electron as a summation over each valley. Utilizing
\eqref{definition:bloch-basis:wavefunction}  and
\eqref{definition:bloch-basis:wavefunction:single-valley} we find
\begin{equation}
  \Psi_{\bm{r}} \equiv \braket{\bm{r} | \Psi} 
  =
  \sum_{q} \Psi^{(q)}_{\bm{r}}.
  \label{equation:wavefunction:multi-valley}
\end{equation}
We expect that the best variational form of function $\Psi^{(q)}$ will be
complicated, mixing together both long- and short-range physics.

Because we expect the $A^{(q)}$ to be very sharply peaked at a particular valley
in the conduction band, we can expand $\Psi^{(q)}$ to partially separate the
long- and short-range scales. In particular, we introduce the "envelope
function" as a shifted inverse truncated transform of the occupation function
\begin{equation}
  \Phi^{(q)}_{\bm{r}}
  \equiv
  \frac{1}{\left(2\pi\right)^{\frac{3}{2}}}
  \int_{BZ}\diff\bm{k}\,
  A^{(q)}_{\bm{k}}
  \mathrm{e}^{\im\left(\bm{k}-\bm{k}_{q}\right)\cdot\bm{r}}
  \label{definition:envelope}
\end{equation}
This is approximately the Fourier transform of $A^{(q)}$ centered on the valley,
and becomes a better approximation the more localized and steeply peaked the
$\bm{k}$-space is. Since we are assuming the occupation function is strongly
peaked in $\bm{k}$-space, the corresponding "envelope function" decays slowly in
real space. In terms of this long-scale envelope function, we can write the
real-space valley wavefunction of the bound electron as
\begin{equation}
  \Psi^{(q)}_{\bm{r}}
  =
  \Phi^{(q)}_{\bm{r}}
  \left( 1 + \varphi^{(q)}_{\bm{r}} \right)
  \mathrm{e}^{\im\bm{k}_{q}\cdot\bm{r}}
  u_{\bm{k}_{q}\bm{r}}^{},
  \label{definition:wavefunction:single-valley}
\end{equation}
where $\varphi^{(q)}$ is such that
\begin{equation}
  \varphi^{(q)}_{\bm{r}}
  \equiv
  \frac
  {
    \int_{BZ}\diff\bm{k}\,
    A^{(q)}_{\bm{k}}
    \mathrm{e}^{\im\left(\bm{k}-\bm{k}_{q}\right)\cdot\bm{r}}
    \left(
      u_{\bm{k}\bm{r}}^{}
      -
      u_{\bm{k}_{q}\bm{r}}^{}
    \right)
  }
  {
    \Phi^{(q)}_{\bm{r}}
    u_{\bm{k}_{q}\bm{r}}
  }.
\end{equation}
We expect $\varphi^{(q)}$ to be quite small, with no zeroth-order contribution
to the integral. While we can expect a first-order contribution, the
localization of $A^{(q)}$ around the minima $\bm{k}_{q}$ will keep this
contribution very small.

The exact form of the potential $\hat{U}$ resulting from deviations of the pure
crystalline structure is still an open question. In this work, we will approach
this question of the form of $\hat{U}$ phenomenologically. We structure our
impurity term such that
\begin{equation}
\mathcal{U}
  =
    \sum_{q,q^\prime}
    \int \diff\bm{r}\,
    \mathrm{e}^{\im\left(\bm{k}_{q^{\prime}}-\im\bm{k}_{q}\right)\cdot\bm{r}}
    U_{\bm{r}}^{(qq^{\prime})}
    u_{\bm{k}_{q}\bm{r}}^{*}
    u_{\bm{k}_{q^{\prime}\bm{r}}}^{}
    \tilde{\Phi}_{\bm{r}}^{(q)}
    \Phi_{\bm{r}}^{\left(q^{\prime}\right)}
\end{equation}
where the $\varphi^{(q)}$ contribution to the integral has been folded into
the potential factor of the integrand as
\begin{equation}
U_{\bm{r}}^{(qq^{\prime})}
  \equiv
    U_{\bm{r}}
    \left( 1 + \varphi_{\bm{r}}^{(q)} \right)^{*}
    \left( 1 + \varphi_{\bm{r}}^{(q^{\prime})} \right).
\end{equation}
This form introduces the idea that the form of the potential may differ for
different combinations of the minima indexed by $q$ and $q^\prime$. In this
work, however, our approach only considers a single effective potential for all
minima combinations. Further work may explore this refinement.

The envelope function also emerges as a convenient way to convert the $k$-space
integral $\mathcal{T}$ to a real-space integral
\begin{equation}
  \mathcal{T}
  =
  \sum_{q,q^\prime}
  \int \diff\bm{r}\,
  \mathrm{e}^{\im\left(\bm{k}_{q^{\prime}}-\bm{k}_{q}\right)\cdot\bm{r}}
  \left\{ \Phi^{(q)}_{\bm{r}} \right\}^{*}
  \left(T\Phi\right)^{(q^{\prime})}_{\bm{r}},
\end{equation}
where we can include the effect of the $E^{(q)}_{\bm{k}}$ band energy as
\begin{equation}
\left(T\Phi\right)^{(q)}_{\bm{r}}
  \equiv
    \frac{1}{\left(2\pi\right)^{\frac{3}{2}}}
    \int_{BZ}\diff\bm{k}\,
    E^{(q)}_{\bm{k}}
    A^{(q)}_{\bm{k}}
    \mathrm{e}^{\im\left(\bm{k}-\bm{k}_{q}\right)\cdot\bm{r}}
\label{expression:kinetic:left}
.
\end{equation}
As we are presuming that the occupation function $A^{(q)}_{\bm{k}}$ is sharply
peaked about the minima of $E^{(q)}_{\bm{k}}$, we can expand the energy around
this minima to lowest order, yielding the following second order expression
\begin{equation}
E^{(q)}_{\bm{k}}
  \approx
    \frac{\hbar^2}{2}
    \left(\bm{k}-\bm{k}_{q}\right)\cdot
    \bm{m}_{q}^{-1}
    \cdot\left(\bm{k}-\bm{k}_{q}\right).
\end{equation}
Inserting this expansion into Eq. \eqref{expression:kinetic:left} we find that
\begin{align}
  \left(T\Phi\right)^{(q)}_{\bm{r}}
  = &
  \frac{1}{\left(2\pi\right)^{\frac{3}{2}}}
  \int_{BZ}\diff\bm{k}\,
  \left\{
    \frac{\hbar^2}{2}
    \left(\bm{k}-\bm{k}_{q}\right)\cdot
    \bm{m}_{q}^{-1}
    \cdot\left(\bm{k}-\bm{k}_{q}\right)
  \right\}
  A^{(q)}_{\bm{k}}
  \mathrm{e}^{\im\left(\bm{k}-\bm{k}_{q}\right)\cdot\bm{r}}
  \notag
  \\
  = &
  \frac{1}{\left(2\pi\right)^{\frac{3}{2}}}
  \int_{BZ}\diff\bm{k}\,
  \left\{
    -\frac{\hbar^2}{2}
    \nabla_{\bm{r}}\cdot \bm{m}_{q}^{-1} \cdot\nabla_{\bm{r}}
  \right\}
  A^{(q)}_{\bm{k}}
  \mathrm{e}^{\im\left(\bm{k}-\bm{k}_{q}\right)\cdot\bm{r}}
  \nonumber
  \\
  = &
  \left\{
    -\frac{\hbar^2}{2}
    \nabla_{\bm{r}}\cdot \bm{m}_{q}^{-1} \cdot\nabla_{\bm{r}}
  \right\}
  \Phi^{(q)}_{\bm{r}}
  \label{opr_func}
\end{align}
This expression is analogous to a kinetic energy operator acting upon the
envelope function, with the quirk that the isotropic inertia mass is replaced
with an effective mass tensor. These effective masses are a result of the
curvature of the conduction band at the minima. We will further explore this
effective mass concept in the context of silicon in the next section.

All together, then, we can restate the above examination succinctly. To wit, our
model of a single electron bound to a "shallow" defect in a crystal is
encompassed in the following equations for the variational wavefunction $\Psi$
and energy $\mathcal{E}$.
\begin{align}
  \Psi_{\bm{r}}
  \approx &
  \sum_{q}
  \Phi^{(q)}_{\bm{r}}
  \mathrm{e}^{\im\bm{k}_{q}\cdot\bm{r}}
  u_{\bm{k}_{q}\bm{r}}^{},
  \label{definition:wavefunction:full}
  \\
  \mathcal{E}
  = &
  \mathcal{T} + \mathcal{U},
  \label{definition:ExpectationEnergy}
  \\
  \mathcal{T}
  \equiv &
  \int \diff\bm{r}\,
  \sum_{q}\sum_{q^{\prime}}
  \mathrm{e}^{\im\left(\bm{k}_{q^{\prime}}-\bm{k}_{q}\right)\cdot\bm{r}}
  \tilde{\Phi^{(q)}}_{\bm{r}}
  \hat{T}^{(q)}
  \left\{ \Phi^{(q^{\prime})}_{\bm{r}}\right\},
  \label{element:kinetic}
  \\
  \hat{T}^{(q)}
  \equiv &
  -\frac{\hbar^2}{2}
  \nabla_{\bm{r}}\cdot \bm{m}_{q}^{-1} \cdot\nabla_{\bm{r}},
  \label{operator:kinetic}
  \\
  \mathcal{U}
  \equiv &
  \int \diff\bm{r}\,
  \sum_{q}\sum_{q^{\prime}}
  \mathrm{e}^{\im\left(\bm{k}_{q^{\prime}}-\bm{k}_{q}\right)\cdot\bm{r}}
  u_{\bm{k}_{q}\bm{r}}^{*}u_{\bm{k}_{q^{\prime}\bm{r}}}^{}
  U_{\bm{r}}
  \tilde{\Phi}_{\bm{r}}^{(q)}
  \Phi_{\bm{r}}^{\left(q^{\prime}\right)}.
\end{align}
In the above, we implicitly assumed that as we are expecting our wavefunction to
be square integrable, it is normalizable and so $\braket{\Psi | \Psi} = 1$. We
also understand that $U_{\bm{r}}$ is not an exact, direct ab-initio expression
of the potential introduced by the crystal defect. Instead we consider a
phenomenological effective potential informed both by ab-initio theoretical
considerations as well as experimental observations.